% casimir_koide_plb.tex -- standalone letter for PLB
\documentclass[final,3p,times,twocolumn]{elsarticle}
\usepackage[T1]{fontenc}
\usepackage{lmodern}
\usepackage{amsmath,amssymb,amsthm,booktabs}
\usepackage[colorlinks=true,citecolor=blue,linkcolor=blue,urlcolor=blue]{hyperref}
\journal{Physics Letters B}

\newcommand{\SU}{\mathrm{SU}}
\newcommand{\UU}{\mathrm{U}}
\newcommand{\Tr}{\mathrm{Tr}}
\newcommand{\adj}{\operatorname{adj}}
\newcommand{\diag}{\mathrm{diag}}
\newcommand{\GeV}{\,\mathrm{GeV}}
\newcommand{\MeV}{\,\mathrm{MeV}}
\newcommand{\Nc}{N_{\!c}}
\newcommand{\Nf}{N_{\!f}}
\newtheorem{theorem}{Theorem}

\begin{document}

\begin{frontmatter}

\title{The Koide ratio as a Casimir invariant:\\
a Diophantine selection of three families}

\author{A.~Rivero}
\address{Departamento de F\'{\i}sica Te\'orica, Universidad de
Zaragoza, 50009 Zaragoza, Spain}
\ead{rivero@unizar.es}

\begin{abstract}
The Koide ratio $K = \sum m_i / (\sum\sqrt{m_i})^2$
equals $2/3$ for charged leptons to $0.001\%$.
We show that this value is the mean-square weight
$\langle J_3^2\rangle$ of the principal $\mathfrak{sl}(2)$
subalgebra of $\mathfrak{sl}(N)$ evaluated in the fundamental
representation, and that the equality
$K_N = \langle J_3^2\rangle$ holds if and only if $N = 3$.
The condition reduces to the Diophantine equation
$N(N^2 - 1) = 24$, whose unique positive integer solution
selects three families.
The operator connecting the two natural mass branches of an
$\SU(3)$ confining theory---the direct ($\adj(M) \propto m$) and
the inverse ($M \propto 1/m$)---is identified as the spin-2
component of $\mathrm{End}(V)$ under the same principal embedding,
linking the mass structure to Kostant's
exponents of $\mathfrak{sl}(3)$.
\end{abstract}

\begin{keyword}
Koide formula \sep family replication \sep
principal $\mathfrak{sl}(2)$ \sep Casimir invariant
\end{keyword}

\end{frontmatter}


%======================================================================
\section{Introduction}
%======================================================================

The charged-lepton mass ratio
\begin{equation}
  K \;\equiv\; \frac{m_e + m_\mu + m_\tau}
  {\bigl(\sqrt{m_e} + \sqrt{m_\mu} + \sqrt{m_\tau}\bigr)^2}
  \;=\; 0.666661
  \label{eq:Kdef}
\end{equation}
was noted by Koide~\cite{Koide:1983} to lie within $0.001\%$
of the value $2/3$.
The precision is striking: among all triplets drawn from
a random mass spectrum, the probability of this level
of agreement is $p < 10^{-5}$~\cite{Rivero:2024}.
Numerous models have been proposed to explain the
relation~\cite{Koide:1983,Sumino:2009,Rivero:2024},
but none has derived the value $2/3$ from first principles
without introducing it as an input.

In this Letter we observe that $2/3$ is a
representation-theoretic invariant: it is the
\emph{mean-square weight} of the fundamental representation
of $\SU(N)$ under its principal $\mathfrak{sl}(2)$
subalgebra~\cite{Kostant:1959}, and that the equality with
the Koide ratio holds \emph{if and only if} $N = 3$.
This provides a Diophantine selection of three families that
is independent of any dynamical model.

We also show that the operator connecting the direct and
inverse mass branches of an $\SU(3)$ confining
theory~\cite{Seiberg:1995,Rivero:2024} is the \emph{spin-2}
component of $\mathrm{End}(V)$ under the principal embedding,
establishing a connection between the Koide mass structure and
Kostant's exponents of $\mathfrak{sl}(3)$.


%======================================================================
\section{The Koide ratio for $\UU(N)$}
\label{sec:KN}
%======================================================================

Let $\hat{z}_i = z_0 + z_i$ ($i = 1, \ldots, N$) be the signed
charges of a $\UU(N)$ family symmetry, with $z_0$ the
$\UU(1)$ generator and $z_i$ the $\SU(N)$ Cartan components.
Canonical normalisation gives:
\begin{equation}
  \sum_{i=1}^{N} z_i = 0\,,\qquad
  \sum_{i=1}^{N} z_i^2 = \frac{1}{2}\,,\qquad
  z_0 = \frac{1}{\sqrt{2N}}\,.
  \label{eq:u3_norm}
\end{equation}

\begin{theorem}[Generalised Koide ratio]
\label{thm:KN}
For the monomial mass law $m_i = \mu\,\hat{z}_i^2$ and the
ratio
\begin{equation}
  K_N \;\equiv\;
  \frac{\sum_{i=1}^N \hat{z}_i^2}
  {\bigl(\sum_{i=1}^N \hat{z}_i\bigr)^2}\,,
\end{equation}
one has $K_N = 2/N$, independently of the Cartan direction.
\end{theorem}

\begin{proof}
Using~\eqref{eq:u3_norm}:
$\sum \hat{z}_i^2 = N z_0^2 + \sum z_i^2 = \tfrac{1}{2} +
\tfrac{1}{2} = 1$
and $\sum \hat{z}_i = N z_0 = \sqrt{N/2}$,
so $K_N = 1/(N/2) = 2/N$.
\end{proof}

For $N = 3$: $K_3 = 2/3$, recovering the Koide value as
an algebraic identity of $\UU(3)$ normalisation.
This is the $n = 2$ uniqueness
theorem~\cite{Koide:1983,Rivero:2024}: among monomial mass
functions $m \propto \hat{z}^n$, only $n = 2$ gives a
direction-independent ratio.


%======================================================================
\section{The principal $\mathfrak{sl}(2)$ and its Casimir}
\label{sec:sl2}
%======================================================================

Every complex simple Lie algebra $\mathfrak{g}$ contains a
distinguished three-dimensional simple subalgebra, the
\emph{principal} $\mathfrak{sl}(2)$, defined (up to conjugacy)
as the unique $\mathfrak{sl}(2)$ for which the fundamental
representation of $\mathfrak{g}$ is
irreducible~\cite{Kostant:1959}.

For $\mathfrak{g} = \mathfrak{sl}(N)$, the fundamental
$\mathbf{N}$ is the spin-$j$ representation of the principal
$\mathfrak{sl}(2)$ with
\begin{equation}
  j = \frac{N-1}{2}\,,\qquad
  \dim = 2j + 1 = N\,.
\end{equation}
The weight operator $J_3$ has eigenvalues
$m = j, j{-}1, \ldots, -j$.
The quadratic Casimir is $C_2 = j(j+1)$, and the
mean-square weight is
\begin{equation}
  \langle J_3^2 \rangle
  \;\equiv\; \frac{\Tr(J_3^2)}{N}
  \;=\; \frac{1}{N}\sum_{m=-j}^{j} m^2
  \;=\; \frac{j(j+1)(2j+1)}{3N}\,.
  \label{eq:J3sq}
\end{equation}
Using $j = (N{-}1)/2$ and $2j{+}1 = N$:
\begin{equation}
  \langle J_3^2 \rangle
  = \frac{(N{-}1)(N{+}1)}{12}
  = \frac{N^2 - 1}{12}\,.
  \label{eq:J3sq_N}
\end{equation}

For $N = 3$: $\langle J_3^2 \rangle = 8/12 = 2/3$.


%======================================================================
\section{The Diophantine condition}
\label{sec:diophantine}
%======================================================================

Setting $K_N = \langle J_3^2 \rangle$:
\begin{equation}
  \frac{2}{N} = \frac{N^2 - 1}{12}
  \quad\Longleftrightarrow\quad
  N(N^2 - 1) = 24\,.
  \label{eq:diophantine}
\end{equation}
Factoring:
\begin{equation}
  N(N-1)(N+1) = 24 = 2 \times 3 \times 4\,.
\end{equation}
The left-hand side is the product of three consecutive integers
centred at $N$.
The only positive integer solution is
$\boxed{N = 3}$
(since $2 \times 3 \times 4 = 24$ and $3 \times 4 \times 5 = 60
> 24$).

Equivalently, the factored form
$(N - 3)(N^2 + 3N + 8) = 0$
has discriminant $\Delta = 9 - 32 < 0$ for the quadratic factor,
so there are no other real roots.

\paragraph{Physical interpretation.}
At $N = 3$ and \emph{only} $N = 3$, the representation dimension
$2j + 1 = N$ coincides with the algebra dimension
$\dim\,\mathfrak{sl}(2) = 3$.
The Casimir per state ($C_2/N$) then equals the Casimir per
generator ($C_2/3$), giving
\begin{equation}
  \frac{C_2}{N} = \frac{C_2}{3}
  = \langle J_3^2 \rangle = K = \frac{2}{3}\,.
\end{equation}
This is the representation-theoretic content of the Koide
value: $K = 2/3$ is the statement that the fundamental
representation of $\SU(3)$ \emph{is} the adjoint of its
principal $\mathfrak{sl}(2)$.


%======================================================================
\section{The spin-2 operator and the mass branches}
\label{sec:spin2}
%======================================================================

In $\SU(3)$ SQCD with $\Nf = \Nc = 3$, the quantum-modified
constraint $\det M = \Lambda^6$~\cite{Seiberg:1995} produces
two holomorphic mass branches from one meson
matrix~\cite{Rivero:2024}:
\begin{equation}
  \langle M_i \rangle \propto \frac{1}{m_i}\,,\qquad
  \adj(M)_i \propto m_i\,.
  \label{eq:branches}
\end{equation}
Define the mixing operator $R$ by
\begin{equation}
  R_{ij} = \langle M_k \rangle
  \quad (k \neq i, j;\; i \neq j)\,,\qquad
  R_{ii} = 0\,.
  \label{eq:Rdef}
\end{equation}
This operator satisfies the exact identity
\begin{equation}
  \boxed{R \cdot \mathbf{M} = 2\,\adj(\mathbf{M})\,,}
  \label{eq:RM}
\end{equation}
where $\mathbf{M} = (M_1, M_2, M_3)^T$ is the meson VEV vector.
On the constraint surface, the Cayley--Hamilton identity gives:
\begin{equation}
  R^3 = \Tr(D^2)\,R + 2\Lambda^6\,I\,,
  \label{eq:cubic}
\end{equation}
where $D = \diag(M_1, M_2, M_3)$.
After conjugation by $D$:
\begin{equation}
  D\,R\,D = \Lambda^6\,(J - I)\,,
  \label{eq:DRD}
\end{equation}
where $J$ is the $3 \times 3$ all-ones matrix.
The right-hand side is the \emph{democratic matrix},
independent of the mass hierarchy: the individual masses
are entirely absorbed by the $D$ conjugation.

\paragraph{$R$ as a spin-2 operator.}
Under the principal $\mathfrak{sl}(2)$ of
Section~\ref{sec:sl2}, the endomorphism algebra
$\mathrm{End}(V)$ for $V$ = spin-1 decomposes as
\begin{equation}
  \mathrm{End}(\text{spin-1})
  = \text{spin-0} \oplus \text{spin-1} \oplus \text{spin-2}\,.
\end{equation}
The three components correspond to:
\begin{itemize}
  \item \textbf{Spin-0} (dim 1): $\propto I$ (trace part).
    Since $\Tr R = 0$, this is absent.
  \item \textbf{Spin-1} (dim 3): antisymmetric matrices.
    Since $R = R^T$, this is absent.
  \item \textbf{Spin-2} (dim 5): symmetric traceless matrices.
    \textbf{This is $R$.}
\end{itemize}
Therefore $R$ is the \emph{unique} spin-2 component of
$\mathrm{End}(V)$ under the principal embedding.

In Kostant's theory~\cite{Kostant:1959}, the adjoint
representation of $\mathfrak{sl}(3)$ decomposes under the
principal $\mathfrak{sl}(2)$ as $\mathbf{8} = \mathbf{3}
\oplus \mathbf{5}$, with the spin content determined by the
exponents $\{m_1, m_2\} = \{1, 2\}$ of $\mathfrak{sl}(3)$.
The spin-1 piece ($\mathbf{3}$) \emph{is} the principal
$\mathfrak{sl}(2)$ itself; the spin-2 piece ($\mathbf{5}$) is
the complementary subspace.
The operator $R$ lives in this exponent-2 representation:
it is the ``second exponent'' of $\mathfrak{sl}(3)$.


%======================================================================
\section{Comparison with other $N = 3$ criteria}
\label{sec:comparison}
%======================================================================

The Casimir--Koide criterion joins four previously known
conditions that each uniquely select $N = 3$:

\begin{center}
\begin{tabular}{lll}
\toprule
Criterion & Equation & Solution \\
\midrule
Diophantine~\cite{Rivero:2024}
  & $rs = 2N$, $r(r{+}1)/2 = 2N$ & $N = 3$ \\
Discriminant
  & $(N{+}2)^2 - 4 = N(2N{+}1)$ & $N = 3$ \\
Binomial
  & $\binom{3N}{N} = (N{+}1) \cdot N(2N{+}1)$ & $N = 3$ \\
CP-phase
  & $N{+}4 = 2N{+}1$ & $N = 3$ \\
\textbf{Casimir--Koide}
  & $\mathbf{N(N^2{-}1) = 24}$ & $\mathbf{N = 3}$ \\
\bottomrule
\end{tabular}
\end{center}

The first four arise from compositeness counting, algebraic
geometry, 11d SUGRA, and the CKM phase, respectively.
The Casimir--Koide criterion is the most algebraically natural:
it identifies the Koide value $2/3$ with a standard invariant of
the principal $\mathfrak{sl}(2)$ embedding, and the $N = 3$
selection follows from the coincidence of the representation and
algebra dimensions.


%======================================================================
\section{Discussion}
\label{sec:discussion}
%======================================================================

The result presented here is purely algebraic: it connects
the Koide ratio (an empirical mass relation) to the Casimir
of the principal $\mathfrak{sl}(2)$ (a canonical object in
Lie theory) through a Diophantine equation that selects $N = 3$.
No dynamical model is assumed beyond the
$\UU(N)$ charge structure and the quadratic mass law
$m_i = \mu\,\hat{z}_i^2$.

The spin-2 identification of the branch-mixing operator $R$
adds a structural layer: in the confining $\SU(3)$ theory, the
operator that connects the direct (lepton-like) and inverse
(quark-like) mass branches is not an arbitrary matrix but the
unique spin-2 irreducible component of $\mathrm{End}(V)$ under
Kostant's principal embedding.
The democratic structure $DRD = \Lambda^6(J - I)$ shows that
the $R$ operator, when expressed in the canonical basis set by
the meson VEVs, becomes mass-independent---the hierarchy is
entirely encoded in the conjugation by $D$.

Two open questions deserve mention.
First, the spin-2 nature of $R$ is suggestive of a gravitational
connection: in 11-dimensional supergravity, the metric $g_{MN}$
is a spin-2 field, while the 3-form $C_3$ is spin-1.
The composites of the $\SU(3)$ theory (mesons coupling to $C_3$,
baryons coupling to $C_6$) are spin-1 objects in the M-theory
lift~\cite{HananyWitten:1997,Rivero:2024}; the branch-mixing
operator $R$ is spin-2, matching the graviton sector.
Whether this is a structural connection or a numerical
coincidence remains to be determined.

Second, the cubic algebra~\eqref{eq:cubic} differs from the
quadratic relation $J^2 = -1$ that characterises $\SU(2)_R$ in
$\mathcal{N} = 2$ supersymmetry.
The emergent structure is therefore not standard $\mathcal{N} = 2$
but a cubic deformation controlled by the mass hierarchy
($\Tr D^2$) and the quantum constraint ($\Lambda^6$).
Whether this cubic defines a consistent deformed superalgebra is
an open problem.


%======================================================================
\section*{Acknowledgments}
%======================================================================

The author thanks M.~Porter for sustained theoretical discussions
over 2011--2024 that identified Seiberg duality and the mesino
interpretation as the mechanism underlying the sBootstrap,
and C.~Brannen for the circulant mass matrix parametrisation.
During the preparation of this work the author used Claude
(Anthropic) for algebraic verification and numerical checks.
The author reviewed and edited all content and takes full
responsibility for the publication.


%======================================================================
\begin{thebibliography}{99}

\bibitem{Koide:1983}
  Y.~Koide,
  Phys.\ Lett.\ B \textbf{120}, 161 (1983).

\bibitem{Sumino:2009}
  Y.~Sumino,
  JHEP \textbf{05}, 075 (2009),
  arXiv:0812.2103.

\bibitem{Rivero:2024}
  A.~Rivero,
  Eur.\ Phys.\ J.\ C \textbf{84}, 1058 (2024),
  arXiv:2407.05397.

\bibitem{Kostant:1959}
  B.~Kostant,
  Amer.\ J.\ Math.\ \textbf{81}, 973 (1959).

\bibitem{Seiberg:1995}
  N.~Seiberg,
  Nucl.\ Phys.\ B \textbf{435}, 129 (1995),
  arXiv:hep-ph/9411149.

\bibitem{HananyWitten:1997}
  A.~Hanany and E.~Witten,
  Nucl.\ Phys.\ B \textbf{492}, 152 (1997),
  arXiv:hep-th/9611230.

\end{thebibliography}

\end{document}
